\section{Degree on a conjugation-invariant part of a word fiber}
\label{sec:degree}

The comparison argument will need a solution that remains in a specified
part of a word map's identity fiber as its coefficients vary. A global
degree calculation does not locate that solution. When the exponent-sum matrix is nonsingular, we prove nonvanishing
for any conjugation-invariant isolating part that contains the identity.

Write $H=\SU(2)$ as the group of unit quaternions. We orient $H$ once and give $H^m$ the product orientation. A \emph{word map} $f:H^m\to H^m$ has coordinates that are words in its $m$ arguments and their inverses, without coefficients. Its exponent-sum matrix $B=(B_{ij})$ records the total exponent of the $j$th argument in the $i$th word.

We use ordinary mapping degree on an isolating open set. Thus, if $M$ is a closed oriented $n$-manifold, $f:M\to M$ is smooth, and $y\notin f(\partial\mathcal O)$, then $\deg(f,\mathcal O,y)$ counts the part of $f^{-1}(y)$ in $\mathcal O$. It is defined even when $\partial\mathcal O$ is not smooth. For the compact selected fiber $K=f^{-1}(y)\cap\mathcal O$, map the local orientation class in $H_n(\mathcal O,\mathcal O\setminus K;\Z)$ to $H_n(M,M\setminus\{y\};\Z)\cong\Z$. This relative-homology construction gives excision and homotopy invariance while the boundary avoids $y$; when $y$ is regular it gives the signed count of preimages. We use the relative orientation class over a compact subset from \cite[Lemma~3.27(a)]{Hatcher2002}, together with excision. Regular fibers and the signed-degree and homotopy-invariance facts are supplied by \cite[Section~2, Lemmas~1--2; Section~5, Theorems~A--B]{Milnor1965}.

The global degree formula for compact connected Lie groups goes back to Gerstenhaber and Rothaus \cite[Theorem~1]{GR1962}; for $H=\SU(2)$ it gives $\deg f=\det B$. A coefficient-free formulation also appears in \cite[Proposition~2.1]{Froyshov2013}. We need a statement about an individual invariant part of the fiber. The following proof provides that localization directly.

\begin{samepage}
\begin{lemma}[Localized word-map degree]
\label{lem:localized-degree}
Let $m\ge1$, and let $f:H^m\to H^m$ be a word map with $\det B\ne0$. Suppose that $\mathcal O\subset H^m$ is open, invariant under simultaneous conjugation, contains $\one$, and satisfies
\[
 \partial\mathcal O\cap f^{-1}(\one)=\varnothing.
\]
Then $\deg(f,\mathcal O,\one)\ne0$.
\end{lemma}
\end{samepage}

\begin{proof}
Let $J=\{e^{i\theta}:\theta\in\R\}$ be the standard maximal torus. On $J^m$, the map $f$ is the torus homomorphism with matrix $B$. In particular, its fiber over $\one$ consists of $\abs{\det B}$ points, each regular for the torus restriction, each with tangent sign $\operatorname{sign}(\det B)$. Let $b$ be the number of these points in $\mathcal O$. The identity is one of them, so $b\ge1$.

Fix an odd prime $p>b$. Simultaneous conjugation by $e^{2\pi i/p}$ generates a cyclic group $C_p$ acting on $H^m$. Its fixed set is $J^m$, and the action is free off $J^m$: since $p$ is prime, a nontrivial stabilizer is the whole group. The word map is equivariant. We shall perturb it equivariantly, without changing its restriction to $J^m$, until $\one$ is regular. The fixed preimages will each contribute $\operatorname{sign}(\det B)$, while the others will contribute in multiples of $p$.

\smallskip
\noindent\emph{Regularizing the fixed preimages.}
At a torus root $x$, use left-translated exponential coordinates in the source and exponential coordinates at $\one$ in the target. Since $x$ commutes with $J$, the coordinate actions are the same adjoint representation:
\[
 \R^m\oplus\C^m,
 \qquad (u,v)\longmapsto (u,e^{4\pi i/p}v).
\]
Here the real directions are the $i$-directions, and each complex plane is spanned by $j,k$. An equivariant derivative has no mixed blocks, because the normal rotation has no fixed vector. Its normal block $N$ is complex linear: commuting with the nonreal scalar $e^{4\pi i/p}$ is equivalent to commuting with multiplication by $i$. The tangent block is $B$.

Choose mutually disjoint invariant coordinate neighborhoods of all torus roots. Each can be chosen entirely inside $\mathcal O$ or entirely outside its closure, because no root lies on the boundary. In one such chart add
\[
 (u,v)\longmapsto s\chi(u,v)(0,v)
\]
to the coordinate expression of $f$, where $\chi$ is an invariant smooth bump function equal to one near the root. Choose $s\in\R$ arbitrarily small with $\det_{\C}(N+sI)\ne0$; only finitely many real values are excluded. This changes neither the torus restriction nor the tangent block. The resulting derivative is invertible, and its sign is
\[
 \operatorname{sign}(\det B)\,
 \operatorname{sign}\det_{\R}(N+sI)
 =\operatorname{sign}(\det B),
\]
since the real determinant of an invertible complex linear map is positive. The coordinate orientations can be chosen compatibly in source and target. Each fixed preimage is now isolated and regular.

\smallskip
\noindent\emph{Regularizing the remaining preimages.}
Choose an invariant open neighborhood of each regular fixed root, with closure contained in a coordinate neighborhood where that root is the only zero. Keep the map unchanged near the closures of the chosen neighborhoods. The zero set outside their union is closed in $H^m$, hence compact. It lies in the free-action locus: a sequence of these remaining zeros approaching $J^m$ would limit to a torus root, already inside one of the chosen neighborhoods.

Cover this compact set by finitely many smaller coordinate neighborhoods, each inside a larger neighborhood whose $p$ translates are disjoint. Different families of translates need not be disjoint from one another. Choose the closures of the bump supports to map into a smaller target logarithm ball, and keep the parameters small enough that every perturbed value stays in the larger exponential coordinate ball. In these target coordinates, add a vector parameter times a bump function in each larger neighborhood, and propagate the change to its translates by equivariance. Choose the bumps positive on the smaller neighborhoods. Near every remaining zero at least one parameter block therefore gives a surjective derivative to the target. By compactness, sufficiently small perturbations introduce no zeros elsewhere outside the protected fixed-root neighborhoods.

The joint map of source point and parameters is consequently transverse to $\one$ near all its zeros. The regular-value theorem and Sard's theorem, applied to the projection of its zero manifold onto parameter space, give arbitrarily small parameter values for which the slice has only regular zeros. These are standard finite-dimensional results; no equivariant transversality theorem is needed on the fixed set. All perturbations can be kept uniformly small enough that $\partial\mathcal O$ stays disjoint from the fiber over $\one$, since this boundary is compact. Homotopy invariance preserves the selected degree.

Every nonfixed zero of the resulting equivariant map occurs in an orbit of $p$ points. The local signs in an orbit agree: conjugation preserves the orientations of both source and target. The signed-count formula thus gives
\[
 \deg(f,\mathcal O,\one)
 \equiv b\,\operatorname{sign}(\det B)\pmod p.
\]
The right side is nonzero modulo $p$, since $0<b<p$. The degree is therefore nonzero.
\end{proof}

\begin{remark}
The same argument works for arbitrarily large odd primes, so it actually gives
$\deg(f,\mathcal O,\one)=b\,\operatorname{sign}(\det B)$.
Only nonvanishing is needed below. When there are no variables and no equations, the corresponding map is the identity of a point, with degree one.
\end{remark}
